\section{Plan del proyecto}

\subsection{Roles y responsabilidades}

\begin{table}[H]
\centering
\caption{Roles del equipo y entregables de cada uno}
\label{tab:roles}
\small
\begin{tabularx}{\linewidth}{@{}lYY@{}}
\toprule
Integrante & Rol y responsabilidad & Entregables \\
\midrule
Zambrano Yong Angel Daniel & Arquitecto y Responsable de Documentación & Diagramas C4, ADR-001 y
ADR-002, modelo de comunicación, CAP/PACELC, justificación distribuida, integración del documento,
tablero y bitácora \\
\addlinespace
Mendoza Bermello Angello Agustín & Líder de Desarrollo & Esqueleto en \texttt{src/}, comunicación
entre servicios, instrucciones de ejecución, capturas \\
\addlinespace
Macías Vásquez César Abraham & Responsable de Calidad & Historias de usuario, RNF con ISO/IEC 25010,
riesgos, trabajos relacionados, pruebas y revisión de coherencia \\
\bottomrule
\end{tabularx}
\end{table}

En un equipo de tres integrantes, el Arquitecto asume además el rol de Responsable de
Documentación.

\subsection{Cronograma de las cuatro entregas}

\begin{table}[H]
\centering
\caption{Cronograma alineado con el calendario del encuadre}
\label{tab:cronograma}
\small
\begin{tabularx}{\linewidth}{@{}llY@{}}
\toprule
Entrega & Semana & Contenido comprometido \\
\midrule
E1 & 5 (23 al 29 de noviembre de 2026) & Definición, requisitos, arquitectura C4, ADR, CAP/PACELC,
esqueleto de dos servicios \\
\addlinespace
E2 & 9 & Microservicios completos, API documentada, contenedores \\
\addlinespace
E3 & 13 & Base de datos distribuida y procesamiento de datos del SLA \\
\addlinespace
E4 & 17 & Sistema final en capas, calidad, pruebas y defensa \\
\bottomrule
\end{tabularx}
\end{table}

\subsection{Tablero de trabajo y tareas}

Las tareas se gestionan como \emph{issues} en GitHub Projects, con un issue por entregable, un
responsable y un criterio de aceptación. Los issues de E1 y de E2 quedan asignados desde la primera
semana, de modo que la segunda entrega no empieza de cero.

\begin{table}[H]
\centering
\caption{Tareas de E1 y E2 y su responsable}
\label{tab:tareas}
\small
\begin{tabularx}{\linewidth}{@{}lYl@{}}
\toprule
Entrega & Tarea & Responsable \\
\midrule
E1 & Estructura del repositorio, README, bitácora y tablero & Angel \\
E1 & Diagrama C4 nivel 1 y ADR-001 & Angel \\
E1 & Diagrama C4 nivel 2, ADR-002 y tabla de comunicación & Angel \\
E1 & Justificación distribuida y análisis CAP/PACELC & Angel \\
E1 & Esqueleto: servicio de técnicos & Angello \\
E1 & Esqueleto: servicio de incidencias y comunicación entre ambos & Angello \\
E1 & Instrucciones de ejecución y capturas & Angello \\
E1 & Historias de usuario con criterios de aceptación & César \\
E1 & Requisitos no funcionales con cifras e ISO/IEC 25010 & César \\
E1 & Análisis de riesgos y trabajos relacionados & César \\
E1 & Revisión de coherencia y pruebas del esqueleto & César \\
E1 & Documento integrado, video y PDF del SGA & Todo el equipo \\
E2 & Contenedores y API Gateway & Angello \\
E2 & Contratos OpenAPI versionados & César \\
E2 & Despliegue en contenedores y pruebas de humo & Angel \\
\bottomrule
\end{tabularx}
\end{table}

\subsection{Reglas de trabajo}

\begin{enumerate}
\item Una rama por tarea y \emph{pull request} a \texttt{main}, aunque la revisión la haga otro
integrante del equipo.
\item Cada integrante confirma con su propia cuenta; ningún integrante sube trabajo de otro.
\item Una entrada en \texttt{BITACORA.md} por sesión, escrita el mismo día.
\item Toda decisión de arquitectura se registra como ADR antes de implementarse.
\item Antes de cada integración: \texttt{git pull --rebase} y compilación del documento.
\end{enumerate}
